% ============================================
% Chapter 6: Conclusion and Reflections
% ============================================

\chapter{Conclusion and Reflections}

\section{Project Summary}

This project designed and implemented NEFT, a simulation platform for electric logistics fleet cooperative dispatching, and evaluated ten dynamic scheduling strategies under controlled, reproducible conditions. The main outcomes are:

\subsection{Simulation System}

\begin{itemize}[leftmargin=4em]
    \item Built a graph-based road network (OpenStreetMap / OSMnx) with Dijkstra shortest-path routing and a simulation-time clock decoupled from wall-clock execution.
    \item Modeled three vehicle types, dynamic task generation with shared seeds, charging-station queuing, and a unified scoring function used for both settlement and batch-selection objectives.
    \item Delivered a snapshot-to-command scheduling interface with shared utility modules, enabling fair comparison and rapid prototyping of new strategies.
    \item Implemented a FastAPI + WebSocket + AMap visualization front-end for real-time monitoring.
\end{itemize}

\subsection{Algorithm Portfolio}

\begin{itemize}[leftmargin=4em]
    \item \textbf{Greedy baselines (4):} Nearest Task First, Heaviest Task First, Earliest Deadline First, Priority First.
    \item \textbf{Search-based batch loaders (3):} DFS Score Search, SA Score Search, RL Batch.
    \item \textbf{Multi-vehicle coordination (3):} Regional Planning, Cluster Auction MAS, Relay Handoff.
    \item All strategies share battery pre-check, mid-route recharge simulation, and charging-station load balancing through common decision templates.
\end{itemize}

\subsection{Experimental Findings}

\begin{itemize}[leftmargin=4em]
    \item Ran 40 headless dynamic benchmarks (4 shared seeds $\times$ 10 algorithms) under identical task sequences per seed.
    \item Dynamic algorithms achieve near-perfect completion; total score and on-time rate separate strategies, especially at medium scale (on-time 41--57\%).
    \item \textbf{DFS Score Search} and \textbf{RL Batch} lead at medium scale ($\approx$3800 points); \textbf{RL Batch} and search variants lead at small scale; \textbf{Cluster Auction MAS} scores highest at large scale on the tested seed.
    \item A supplementary Gurobi static solver baseline did not consistently outperform dynamic policies; at small and medium scales several dynamic strategies scored higher, possibly because time-penalty terms in the MIP model distort the objective relative to the online scoring function.
\end{itemize}

\section{Lessons Learned}

\begin{enumerate}[leftmargin=4em]
    \item \textbf{Fair comparison requires infrastructure, not just algorithms.} The snapshot interface and shared task seeds were essential; without them, score differences would reflect random task noise rather than strategy quality.
    \item \textbf{Batch loading dominates medium-scale performance.} When pending queues grow, explicit search over feasible batches (DFS / RL) substantially beats single-key greedy sorting.
    \item \textbf{Cooperation has a scale threshold.} Relay and auction mechanisms need sufficient fleet size and geographic task spread; otherwise coordination overhead hurts more than it helps.
    \item \textbf{Static optima $\neq$ dynamic benchmark winners.} Exact solvers are only meaningful when the mathematical model faithfully encodes the same penalties and timing semantics as the simulator; mismatched time-penalty modeling can make ``optimal'' static plans score worse than simple online heuristics.
    \item \textbf{Visualization accelerates debugging.} Map-based route display exposed charging and routing edge cases that log inspection alone would have missed.
\end{enumerate}

\section{Future Work}

\begin{itemize}[leftmargin=4em]
    \item Refine the static MIP formulation so that deadline and overdue penalties align exactly with the dynamic scoring module, enabling a fair static-vs-dynamic gap analysis.
    \item Extend benchmarks with more seeds per scale and report mean $\pm$ std.\ dev.\ for statistical confidence.
    \item Tune cooperative strategies (relay thresholds, auction bid weights) and test on denser task streams where handoff should pay off.
    \item Explore hybrid policies that switch between greedy, search, and multi-agent modes based on pending-queue size and fleet utilization.
\end{itemize}

\section{Team Reflections}

This project took us from textbook scheduling concepts to a working system where vehicles move on a real map, charge at stations, and compete under a unified score. Implementing ten strategies on one interface showed how much engineering effort lies beyond the core algorithm --- battery safety, seed reproducibility, and simulation-time semantics all shaped the final rankings as much as the dispatching logic itself.

The experimental phase changed how we evaluate algorithms: a single impressive run is misleading, but shared seeds across ten strategies revealed clear patterns --- search wins at medium load, coordination matters at large scale, and simple nearest-neighbor rules remain a strong baseline. Working in a team mirrored real software practice: simulation core, algorithms, visualization, and analysis had to integrate through clear contracts, much like the snapshot-to-command design we built for the schedulers themselves.
